% 张三慧V2第18章PDF物理页156源思考题18.1--18.11；原图18.18/18.19需复核。
\section*{第18章\quad 磁力：思考题解析}
\paragraph{思考题 18.1}
磁力 \(\boldsymbol F=q\boldsymbol v\times\boldsymbol B\)，对于固定速率与电荷大小，有 \(\max|\boldsymbol F|=|q|v|\boldsymbol B|\)，当 \(\boldsymbol v\perp\boldsymbol B\) 取最大。因此 \(\boxed{B=F_{\max}/(|q|v)}\)，且 \(\boldsymbol B\) 是按右手定则确定方向的矢量。原书写 \(B=F_{\max}\times v/qv^2\) 形式需区分矢量积方向与标量大小。

\paragraph{思考题 18.2}
地球的偶极磁场弯曲带电宇宙射线轨迹，低磁刚度粒子在近赤道区域更容易受磁屏蔽；近磁极的场线更容易与大气相连，某些入射轨迹可沿场进入极区。确切穿透条件取决于地磁截止刚度、入射方向和粒子电荷；并不是所有带电粒子只在两极出现。

\paragraph{思考题 18.3}
在赤道，地球磁场主要近似水平朝北。取东 \(+\hat x\)、北 \(+\hat y\)、上 \(+\hat z\)，正质子沿东飞行时 \(\boldsymbol F=qv\hat x\times B\hat y=+qvB\hat z\)，受力向上；沿西受力向下。若希望洛伦兹磁力\emph{补偿重力}，应沿东发射且使 \(qvB\approx mg\)。但对单个质子重力很弱、磁场轨道和地球弯曲需单独分析，不会真的沿赤道匀速绕地球不变。

\paragraph{思考题 18.4}
仍取东 \(+\hat x\)、北 \(+\hat y\)、上 \(+\hat z\)，电场 \(\boldsymbol E=-E\hat z\)，地磁 \(\boldsymbol B=B\hat y\)。无偏转需 \(\boldsymbol E+\boldsymbol v\times\boldsymbol B=0\)。若向东 \(\boldsymbol v=v\hat x\)，\(\boldsymbol v\times\boldsymbol B=+vB\hat z\)，故 \(\boxed{\boldsymbol v\text{ 向东},\ v=E/B}\)。由于电荷因子 \(q\) 同时乘两种力，正负电子在同一速度下也可无偏转，不能只按电子负电改向西。

\paragraph{思考题 18.5}
理想静磁 Lorentz 力始终与速度垂直，\(\boxed{\boldsymbol F_B\cdot\boldsymbol v=q(\boldsymbol v\times\boldsymbol B)\cdot\boldsymbol v=0}\)，故纯静磁场不能改变粒子动能。射频加速器中提供能量的是时变电场，磁场只弯转或约束轨道；运动粒子在辐射损失后动能甚至可能减小。

\paragraph{思考题 18.6}
若磁场空间变化缓慢，粒子回旋磁矩 \(\boxed{\mu_{\rm ad}=mv_\perp^2/(2B)}\) 近似守恒。粒子由弱B到强B时 \(v_\perp^2\propto B\) 增大，若无电场且能量守恒，\(v_\parallel^2=v^2-v_\perp^2\) 下降，最终可能在 \(v_\parallel=0\) 发生磁镜反射。总 \(K=mv^2/2\) 近似不变。若场剧变或强碰撞，绝热近似失效。

\paragraph{思考题 18.7}
长直导线电流 \(I\) 的场 \(\boldsymbol B=\mu_0I\hat\varphi/(2\pi r)\) 随半径变化且沿圆周。对沿轴向运动的正粒子，\(\boldsymbol v=v_z\hat z\)，磁力 \(q\boldsymbol v\times\boldsymbol B=-qv_z B_\varphi\hat r\)（对电流同向的粒子径向向内）。因为 \(B\propto1/r\)，轨迹并非均匀B中的等半径圆周，若粒子获得径向速度，还会产生其他方向力分量。

\paragraph{思考题 18.8}
在互相垂直均匀场中，为使不偏转，要求 \(\boldsymbol F=q(\boldsymbol E+\boldsymbol v\times\boldsymbol B)=0\)。入射方向须令 \(\boldsymbol v\times\boldsymbol B=-\boldsymbol E\)，从而速率 \(\boxed{v=E/B}\)（\(\boldsymbol v\perp\boldsymbol B\)）。对正负粒子同样适用，选中速度与 \(q,m\) 无关；但若粒子入射方向不合，仍会偏转。

\paragraph{思考题 18.9}
导体总体安培力 \(\boldsymbol F=I\boldsymbol L\times\boldsymbol B\) 的电流 \(I\) 按\emph{正电荷电流方向}定义；电子实际漂移方向与电流相反且电荷为负，两负号抵消，故给定I、B时安培力方向\emph{与载流子的正负无关}。霍尔电压则依赖不同符号载流子的分离方向，单一载流子模型 \(R_H=1/(nq)\) 符号随 \(q\) 改变。

\paragraph{思考题 18.10}
相同静磁方向下，电荷正负对应相反曲率符号，因为 \(\boldsymbol F=q\boldsymbol v\times\boldsymbol B\)。对垂直纸面向外的 \(B\)，需将图18.18上的\emph{运动箭头}与转向结合，按右手定则辨认正电子与电子；不能仅凭某一帧的螺旋“顺时针”而忽略运动方向。电离损失使粒子 \(p_\perp\) 下降，半径 \(r=p_\perp/(|q|B)\) 逐渐缩小，因此常出现向内缠绕的旋线。原图纸面箭头方向仍须逐图核对。

\paragraph{思考题 18.11}
静电场沿 \(+\hat y\)，正质子总机械能 \(\boxed{K+q\phi=\mathrm{const}}\)，其中 \(\phi(y)=-Ey\)，故上升到更高 \(y\) 时势能下降、\(\boxed{K\text{ 增大}}\)；下降时动能减小，\emph{如果E确为向上}。磁力不做功，只改变速度方向。图18.19中沿x方向的尖点回返在交叉场的旋轮线（trochoid/cycloid）轨迹中出现，\(x\) 向速度可改号，而速度大小由 \(y\) 决定。需结合初速度为零的源图设置才可指定每个拐点的绝对速率。
